% Part: first-order-logic
% Chapter: sequent-calculus
% Section: quantifier-rules

\documentclass[../../../include/open-logic-section]{subfiles}

\begin{document}

\olfileid[fr]{fol}{seq}{qrl}

\olsection{Règles pour les quantificateurs}

\subsection{Règles pour $\lforall$}

\begin{defish}
\Axiom$ !A(t), \Gamma \fCenter \Delta$
\RightLabel{\LeftR{\lforall}}
\UnaryInf$ \lforall[x][!A(x)],\Gamma \fCenter \Delta$
\DisplayProof
\hfill
\Axiom$ \Gamma \fCenter \Delta, !A(a) $
\RightLabel{\RightR{\lforall}}
\UnaryInf$ \Gamma \fCenter \Delta, \lforall[x][!A(x)]$
\DisplayProof
\end{defish}

Dans \LeftR{\lforall}, $t$ est un terme clos, c'est-à-dire un terme
sans variable. Dans \RightR{\lforall}, $a$~est !!a{constant} qui ne
doit figurer nulle part dans le séquent inférieur de la règle
\RightR{\lforall}. On appelle $a$ l'\emph{eigenvariable} de
l'inférence \RightR{\forall}.\footnote{On emploie le terme
« eigenvariable » pour des raisons historiques, bien que $a$ soit
!!a{constant} dans la règle ci-dessus.}

\subsection{Règles pour $\lexists$}

\begin{defish}
\Axiom$ !A(a), \Gamma \fCenter \Delta $
\RightLabel{\LeftR{\lexists}}
\UnaryInf$ \lexists[x][!A(x)], \Gamma \fCenter \Delta$
\DisplayProof
\hfill
\Axiom$ \Gamma \fCenter \Delta, !A(t) $
\RightLabel{\RightR{\lexists}}
\UnaryInf$ \Gamma \fCenter \Delta, \lexists[x][!A(x)]$
\DisplayProof
\end{defish}

De même, $t$~est un terme clos et $a$~est !!a{constant} qui ne
figure pas dans le séquent inférieur de la règle \LeftR{\lexists}.
On appelle $a$ l'\emph{eigenvariable} de l'inférence
\LeftR{\lexists}.

La condition selon laquelle une eigenvariable ne doit pas figurer
dans le séquent inférieur de l'inférence \RightR{\lforall} ou
\LeftR{\lexists} s'appelle la \emph{condition sur l'eigenvariable}.

\begin{explain}
Rappelons la convention suivante : lorsque $!A$ est !!a{formula}
dans laquelle la !!{variable}~$x$ est libre, on l'indique en
écrivant~$!A(x)$. Dans ce même contexte, $!A(t)$ abrège alors
$\Subst{!A}{t}{x}$. On pourrait donc aussi écrire la règle
\RightR{\lexists} sous la forme suivante :
\begin{prooftree}
  \Axiom$\Gamma \fCenter \Delta, \Subst{!A}{t}{x}$
  \RightLabel{\RightR{\lexists}}
  \UnaryInf$\Gamma \fCenter \Delta, \lexists[x][!A]$
\end{prooftree}
Remarquons que $t$ peut déjà figurer dans~$!A$ ; par exemple,
$!A$~pourrait être~$\Atom{\Obj P}{t,x}$. Ainsi, déduire
$\Gamma \Sequent \Delta, \lexists[x][\Atom{\Obj P}{t,x}]$
de~$\Gamma \Sequent \Delta, \Atom{\Obj P}{t,t}$ constitue une
application correcte de~\RightR{\lexists} : on peut « remplacer »
une ou plusieurs occurrences de~$t$ dans la prémisse, sans
nécessairement les remplacer toutes, par la !!{variable} liée~$x$.
En revanche, les conditions sur l'eigenvariable dans
\RightR{\lforall} et~\LeftR{\lexists} imposent que la
!!{constant}~$a$ ne figure pas dans~$!A$. On ne peut donc pas déduire
correctement $\Gamma \Sequent \Delta, \lforall[x][\Atom{\Obj P}{a,x}]$
de $\Gamma \Sequent \Delta, \Atom{\Obj P}{a,a}$ à l'aide
de~$\RightR{\lforall}$.
\end{explain}

\begin{explain}
Dans \RightR{\lexists} et \LeftR{\lforall}, aucune autre restriction
ne porte sur le terme clos~$t$. En revanche, dans les règles
\LeftR{\lexists} et \RightR{\lforall}, la condition sur
l'eigenvariable impose en particulier que la !!{constant}~$a$ ne
figure nulle part en dehors de~$!A(a)$ dans le séquent supérieur.
Cette condition est nécessaire à la correction du système : elle
garantit que celui-ci ne !!{derive}s que des séquents valides. Sans
elle, les inférences suivantes seraient permises :
\begin{prooftree}
  \Axiom$!A(a) \fCenter !A(a)$
  \RightLabel{*\LeftR{\lexists}}
  \UnaryInf$\lexists[x][!A(x)] \fCenter !A(a)$
  \RightLabel{\RightR{\lforall}}
  \UnaryInf$\lexists[x][!A(x)] \fCenter \lforall[x][!A(x)]$
  \DisplayProof\bottomAlignProof
  \qquad
  \Axiom$!A(a) \fCenter !A(a)$
  \RightLabel{*\RightR{\lforall}}
  \UnaryInf$!A(a) \fCenter \lforall[x][!A(x)]$
  \RightLabel{\LeftR{\lexists}}
  \UnaryInf$\lexists[x][!A(x)] \fCenter \lforall[x][!A(x)]$
\end{prooftree}
Or le schéma de séquent
$\lexists[x][!A(x)] \Sequent \lforall[x][!A(x)]$ n'est pas valide
en général.
\begin{editorial}
  La source dit ici que le terme n'est soumis à aucune restriction ;
  on rappelle qu'il doit rester clos, conformément aux règles
  précédentes. L'absence de la constante en dehors de la formule
  indiquée est une conséquence de la condition sur l'eigenvariable,
  et non sa formulation complète : elle doit être absente de tout
  le séquent inférieur. Le dernier séquent n'est pas valide en tant
  que schéma ; certaines de ses instances peuvent être valides.
\end{editorial}
\end{explain}

\end{document}
